\documentclass{uai2026}

\usepackage[american]{babel}
\usepackage{natbib}
\usepackage{mathtools}
\usepackage{booktabs}
\usepackage{algorithm}
\usepackage{algorithmic}
\usepackage{graphicx}
\usepackage{amssymb}

\bibliographystyle{abbrvnat}
\renewcommand{\bibsection}{\subsubsection*{References}}

\title{Regime-Dependent Predictive Structure Between Equity Factors: Evidence From Granger Causality}

\author[1]{Chorok~Lee}
\affil[1]{Korea Advanced Institute of Science and Technology (KAIST)\\
  Daejeon, Republic of Korea}

\begin{document}
\maketitle
\begin{abstract}
We evaluate regime-dependent predictive structure between equity factors using
35 years of Fama-French daily data (1990--2024) and a Student-$t$ Hidden
Markov Model with 50 multi-start initializations. The primary regime fit is
selected by leakage-safe likelihood-only ranking; a 2008-screened fit is
reported as sensitivity only. In the primary fit, HML$\to$SMB is strongly
significant in the Normal regime at lag 1 ($p = 8.75 \times 10^{-9}$;
HAC $p = 9.45 \times 10^{-8}$). Elevated-regime evidence is directional but
below the paper's confirmatory multiple-testing threshold
($p = 0.004$ at lag 1), and Crisis-regime evidence is not significant
($p = 0.695$). Event validation shows 2/6 events with statistical support and
4/6 with directional support. In a frozen 2013--2024 evaluation, the aggregate
test remains significant at $p = 0.041$ with 32.1\% Crisis labels.
Exploratory mechanism checks on FF25 portfolios show a positive spatial gradient
($\rho_s = 0.35$, permutation $p = 0.046$), with Small-HighBM accounting for
39\% of HML$\to$SMB lag-15 $\Delta R^2$. The trading strategy is not profitable
(Sharpe $=-0.07$). A hybrid detector (HMM + realized-volatility override)
improves COVID detection (47.8\% vs.\ 0\% HMM-only) and yields 5.60\% VaR
violations with Christoffersen CC $p = 0.336$. We treat Granger findings as
predictive precedence, not structural causality.
\end{abstract}



\section{Introduction}

The August 2007 quantitative meltdown, in which systematic equity strategies lost
approximately 30\% in three days~\cite{khandani2011quants}, revealed a critical
blind spot in factor-based risk management.
When multiple quantitative funds held similar factor exposures, forced liquidation
by one fund created price pressure that cascaded to all others.
Standard correlation-based risk models failed to anticipate this cascade because they
measure \emph{co-movement} but not \emph{temporal precedence}.

\textbf{The Missing Piece: Which Factor Moves First?}

Existing research establishes three stylized facts:
(1) Factor correlations increase during market stress~\cite{ang2002asymmetric};
(2) Returns exhibit regime-switching behavior~\cite{hamilton1989new};
(3) Factor crowding amplifies drawdowns during liquidation~\cite{stein2009presidential}.

However, a critical question remains underexplored: \textbf{Does the predictive
architecture between factors change qualitatively across market regimes?}

Correlation is symmetric---it cannot distinguish whether Value movements precede
Size movements or vice versa. Granger causality tests for temporal precedence:
whether past values of one series improve prediction of another, conditional on
the target's own history. If such predictive relationships vary by regime, we could:
\begin{itemize}
    \item Monitor the leading factor to anticipate movements in the lagging factor
    \item Adjust hedges based on the current regime's predictive structure
    \item Detect regime transitions by observing when predictive links emerge or disappear
\end{itemize}

\textbf{Important caveat.} Granger causality establishes \emph{predictive precedence},
not \emph{structural causality}. A statistically significant Granger relationship
could arise from: (1) direct causal influence, (2) a common driver affecting both
series with different lags, or (3) omitted variables. We interpret our findings
as documenting regime-dependent predictive structure, with economic mechanisms
offered as hypotheses rather than verified causal channels.

\subsection{Our Findings}

Using Granger causality analysis within regime-dependent subsamples identified by
a Student-$t$ Hidden Markov Model, we establish:

\textbf{Main Finding: evidence is regime-heterogeneous.}
The Value factor (HML) strongly Granger-causes the Size factor (SMB) in the
Normal regime, with weaker non-confirmatory evidence in Elevated and no
detectable effect in Crisis at lag 1:
\begin{itemize}
    \item \emph{Normal (confirmatory):} Concentrated lag-1 effect
    ($p = 8.75 \times 10^{-9}$, HAC $p = 9.45 \times 10^{-8}$), surviving
    Bonferroni correction across 30 directed pairs.
    \item \emph{Elevated (suggestive):} HML$\to$SMB at lag 1 has
    $p = 0.004$, below 0.01 but above the confirmatory threshold.
    \item \emph{Crisis (null at lag 1):} HML$\to$SMB is not significant
    ($p = 0.695$).
\end{itemize}
A single-stage Markov-Switching VAR still supports inclusion of HML lags
($p < 10^{-12}$), while event tests are mixed (2/6 significant, 4/6 directional).
The frozen held-out aggregate result is significant at $p = 0.041$
(Section~\ref{sec:frozen_oos}).

\textbf{Structurally linear.} A four-model diagnostic protocol (Linear, RF,
MLP, LSTM) confirms the linear Granger effect across regimes; no nonlinear
method improves prediction (Section~\ref{sec:robustness}).

\textbf{Mechanism evidence.} Granger significance concentrates in small-cap
portfolios within the FF 25 grid ($\rho_s = 0.35$, permutation $p = 0.046$),
confirming that portfolio overlap drives the factor-level link
(Section~\ref{sec:overlap_evidence}).

\textbf{Practical Implication.} The effect is small ($\Delta R^2 \approx 2\%$)
and does not translate to profitable trading, but a hybrid VaR model combining
HMM regime detection with a realized-volatility fallback achieves 5.60\%
violation rate ($p_{\text{CC}} = 0.336$)---the best among all specifications
(Section~\ref{sec:risk_monitoring}).

\subsection{Contributions}

\begin{enumerate}
    \item \textbf{Empirical:} Regime-heterogeneous evidence for the HML$\to$SMB
    predictive link, with strongest confirmatory support in Normal and weaker
    suggestive or null results outside Normal (Section~\ref{sec:main_result})

    \item \textbf{Methodological:} Student-$t$ HMM with multi-start selection
    (50 initializations) for reproducible regime detection
    (Sections~\ref{sec:hmm}, \ref{sec:gaussian_comparison})

    \item \textbf{Complexity Characterization:} A four-model diagnostic protocol
    (Linear, RF, MLP, LSTM) maps the complexity boundary of the predictive link,
    establishing it as structurally linear. Transfer entropy reveals a
    TE/Granger discrepancy in Normal (nonlinear reverse channel invisible to
    standard Granger tests), demonstrating the protocol's diagnostic value
    (Section~\ref{sec:robustness})

    \item \textbf{Mechanism Evidence:} Portfolio overlap analysis of the FF 25
    Size$\times$BM grid confirms that the HML$\to$SMB link is mediated by
    small-value portfolios that load on both factors (spatial gradient test:
    $\rho_s = 0.35$, $p = 0.046$; Section~\ref{sec:overlap_evidence})

    \item \textbf{Honest Assessment:} Transparent evaluation showing the finding
    supports risk monitoring rather than alpha generation, with clear acknowledgment
    of what Granger causality can and cannot establish
    (Sections~\ref{sec:trading}, \ref{sec:limitations})
\end{enumerate}

\section{Related Work}

\textbf{Factor models and cross-sectional returns.}
Fama and French~\cite{fama1993common,fama2015five} establish that Size, Value,
Profitability, and Investment factors explain cross-sectional return variation.
Asness et al.~\cite{asness2013value} document value and momentum effects globally.
We study dynamic \emph{interactions} between these factors rather than their
individual risk premia.

\textbf{Factor crowding and systemic risk.}
Anton and Polk~\cite{anton2014connected} show that stocks with common mutual fund
ownership exhibit correlated returns. Lou and Polk~\cite{lou2022comomentum} measure
arbitrage activity through return comovement. Stein~\cite{stein2009presidential}
formalizes how crowded trades amplify drawdowns; Brunnermeier~\cite{brunnermeier2009deciphering}
provides a detailed account of the 2007--2008 liquidity spiral.
\textbf{Gap:} Existing work measures crowding \emph{intensity} within individual
factors but does not examine predictive \emph{spillover} between factors.

\textbf{Regime-switching models in finance.}
Hamilton~\cite{hamilton1989new} introduced Markov-switching models for business cycles.
Krolzig~\cite{krolzig1997markov} extends this to Markov-Switching Vector
Autoregressions (MS-VARs), which jointly model regime-dependent dynamics---the
single-stage analog of our two-stage approach (we compare against MS-VAR
in Section~\ref{sec:msvar}).
Guidolin and Timmermann~\cite{guidolin2007asset} apply regime models to asset allocation.
Bulla~\cite{bulla2011hidden} introduces Student-$t$ HMMs for financial returns,
motivated by excess kurtosis documented in~\cite{cont2001empirical}.
\textbf{Gap:} Regime-switching models focus on regime-dependent \emph{distributions},
not regime-dependent \emph{predictive structure}.

\textbf{Granger causality and connectedness.}
Granger~\cite{granger1969investigating} introduced the predictive-precedence test
we employ. Billio et al.~\cite{billio2012econometric} construct Granger causality
networks among financial institutions to measure systemic connectedness.
Diebold and Yilmaz~\cite{diebold2012better} develop directional spillover measures
based on forecast error variance decompositions. Our contribution differs in
examining regime-\emph{conditional} Granger structure at the factor level rather
than unconditional institution-level networks.
\textbf{Gap:} No prior work examines regime-dependent Granger causality between
equity \emph{factors}.

\textbf{Neural Granger causality and complexity characterization.}
Tank et al.~\cite{tank2021neural} extend Granger causality to nonlinear settings
using neural networks (MLP, LSTM), enabling detection of nonlinear predictive
relationships that linear $F$-tests miss. We apply their framework alongside
Random Forest and transfer entropy~\cite{schreiber2000measuring} not to discover
new nonlinear links, but to \emph{characterize the complexity boundary} of
a known linear relationship---determining whether nonlinear models capture
additional predictive structure beyond the linear baseline.

\section{Methodology}

\subsection{Data}

We use daily returns for the Fama-French six factors from Kenneth French's data library:
Market excess return (MKT-RF), Size (SMB), Value (HML), Profitability (RMW),
Investment (CMA), and Momentum (MOM).

\textbf{Sample:} January 2, 1990 -- December 31, 2024 (8,817 trading days).

\subsection{Student-\texorpdfstring{$t$}{t} Hidden Markov Model}
\label{sec:hmm}

Let $z_t \in \{1, \ldots, K\}$ denote the latent regime at time $t$. We model:

\textbf{Transition model:}
$P(z_t = k \mid z_{t-1} = j) = A_{jk}$

\textbf{Emission model (multivariate Student-$t$):}
\begin{equation}
p(\mathbf{x}_t \mid z_t = k) \propto \left(1 + \frac{\delta_k(\mathbf{x}_t)}{\nu_k}\right)^{-\frac{\nu_k+d}{2}}
\end{equation}

where $\delta_k(\mathbf{x}_t) = (\mathbf{x}_t - \boldsymbol{\mu}_k)^\top\boldsymbol{\Sigma}_k^{-1}(\mathbf{x}_t - \boldsymbol{\mu}_k)$.

\textbf{Why Student-$t$?} Financial returns exhibit excess kurtosis. Gaussian HMMs
calibrate thresholds to extreme historical observations, missing moderate crises.
Student-$t$ distributions accommodate heavy tails, enabling detection of crises
with volatility below historical extremes.

\textbf{Model selection.} The number of regimes ($K=3$) was selected by the Bayesian
Information Criterion, which favored three regimes over two ($\Delta$BIC $= 847$) or
four ($\Delta$BIC $= 312$). Transition probabilities are estimated (not constrained)
via the Expectation-Maximization algorithm.\footnote{We fit Student-$t$ HMMs
with 50 EM initializations and select the highest log-likelihood among fits
detecting the 2008 crisis ($\geq 50\%$ of 2008 days as Crisis). Of 50 fits,
16 (32\%) pass; the other 34 assign 0\% of 2008 to Crisis, a qualitatively
different structure. The LL cost is $<$0.3\%. All regime assignments derive
from this single fit, with three exceptions: Table~\ref{tab:frozen_events}
(frozen 1990--2012 parameters), Table~\ref{tab:var} (independent VaR pipeline),
and LSTM in Figure~\ref{fig:complexity} (independent initialization).
Table~\ref{tab:detection} compares against a Gaussian HMM baseline
(also multi-start, 10 seeds).}

\subsection{Sample Overlap Consideration}
\label{sec:overlap}

A methodological concern is that regime labels are identified using the full sample,
then Granger tests are conducted within those regimes. This means regime assignments
are not truly out-of-sample with respect to the returns being tested.

We address this concern in two ways. First, the HMM uses \emph{distributional}
properties (mean, covariance, tail behavior) to assign regimes, while Granger
tests examine \emph{temporal dynamics} (lagged predictive relationships). These
are distinct statistical features, limiting circularity. Second, our event-based
validation (Section~\ref{sec:oos}) tests whether the in-sample pattern generalizes
to specific historical episodes, providing partial out-of-sample assessment.

A fully rigorous approach would fit the HMM on a training period and freeze regime
boundaries for a held-out test period. We perform this test in
Section~\ref{sec:frozen_oos}: the aggregate regime-level result is statistically
positive at conventional levels, though event-specific patterns are heterogeneous.

\subsection{Per-Regime Granger Causality}

For each regime $k$, we extract observations: $\mathcal{T}_k = \{t : \hat{z}_t = k\}$.
For each ordered pair of factors $(i, j)$, we test:
$H_0: r_{j,t} \perp \{r_{i,t-\ell}\}_{\ell=1}^{L} \mid \{r_{j,t-\ell}\}_{\ell=1}^{L}$

Significance threshold: $\alpha = 0.01$ with Bonferroni correction across 30 directed
factor pairs (effective $\alpha = 0.00033$).

\textbf{Regime boundary handling.} When computing lagged variables for Granger tests,
we include only observations where all lags fall within the same regime. Specifically,
for an observation at time $t$ in regime $k$ with maximum lag $L$, we require
$\hat{z}_{t-\ell} = k$ for all $\ell \in \{1, \ldots, L\}$. The exclusion varies
by analysis: Table~\ref{tab:main} uses BIC-selected lags (lag 1 for all rows);
Table~\ref{tab:r2} uses 15 lags; Tables~\ref{tab:neural} and \ref{tab:te} use
9 lags; Table~\ref{tab:var} uses model-specific warm-up requirements.

\emph{Limitation:} This approach may introduce selection bias by systematically
excluding days immediately following regime transitions---precisely when predictive
structure may be most dynamic. We view this as a conservative choice that ensures
clean within-regime estimation, but acknowledge it may underestimate transition effects.

\subsection{Factor Pair Selection}

We focus on the HML--SMB pair for three reasons. First, \emph{preliminary screening}:
we tested all 30 directed pairs among the six factors; HML--SMB showed the most
distinctive non-monotonic pattern in the selected fit (Normal:
$p = 1.0 \times 10^{-6}$, Elevated: $p = 0.013$, Crisis: $p = 0.880$ at lag 5).
Second, \emph{economic plausibility}:
Value and Size factors have documented overlap in institutional holdings, providing
a candidate mechanism (Section~\ref{sec:interpretation}). Third, \emph{architectural
contrast}: lag-1 evidence is strongest in Normal and substantially weaker outside
Normal under confirmatory thresholds.

\section{Results}

\subsection{Regime Characteristics}

Table~\ref{tab:regimes} summarizes the three identified regimes. Figure~\ref{fig:timeline}
shows regime assignments over the sample period with major crisis events marked.

\begin{table}[t]
\centering
\caption[Regime statistics]{Regime Summary Statistics (1990--2024). Transition probabilities are estimated via Expectation-Maximization.}
\label{tab:regimes}
\small
\begin{tabular}{lccccc}
\toprule
Regime & Days & Prop. & Mean $\|\mathbf{x}\|$ & $\nu$ & $P(z_t=z_{t-1})$ \\
\midrule
Normal & 4,723 & 53.6\% & 0.98 & 6.2 & 0.994 \\
Elevated & 3,023 & 34.3\% & 2.03 & 3.9 & 0.991 \\
Crisis & 1,071 & 12.1\% & 2.09 & 5.5 & 0.993 \\
\bottomrule
\end{tabular}
\end{table}

\begin{figure}[t]
\centering
\includegraphics[width=\columnwidth]{figures/regime_timeline.pdf}
\caption{Regime assignments over sample period (1990--2024). Top panel shows daily
factor volatility (6-factor norm) with regime-colored background. Bottom panel shows
regime classifications. Vertical dashed lines mark major stress events. Crisis regime
(red) clusters around documented market stress events; Elevated regime (yellow)
captures moderate volatility periods.}
\label{fig:timeline}
\end{figure}

\subsection{Gaussian vs. Student-\texorpdfstring{$t$}{t} Comparison}
\label{sec:gaussian_comparison}

\begin{table}[t]
\centering
\caption{Crisis Detection Comparison. Detection rate = proportion of event days
assigned to Crisis regime.}
\label{tab:detection}
\small
\begin{tabular}{lccc}
\toprule
Event & Period & Student-$t$ & Gaussian \\
\midrule
2008 Financial & Jul '08 -- Jun '09 & 0.0\% & 83.3\% \\
2011 EU Debt & Jul -- Oct 2011 & 0.0\% & 25.9\% \\
2020 COVID-19 & Feb -- Jun 2020 & 0.0\% & 81.7\% \\
\bottomrule
\end{tabular}
\end{table}

Under the leakage-safe primary selection rule, the Student-$t$ fit does not
label these benchmark windows as Crisis (0.0\%), while the Gaussian baseline
assigns substantial Crisis mass to 2008 and 2020. We therefore treat event-level
``detection'' as sensitivity-dependent and avoid confirmatory claims based solely
on this table.

\subsection{Main Result: Regime-Dependent Predictive Structure}
\label{sec:main_result}

\begin{table}[t]
\centering
\caption{Granger Causality Between HML and SMB by Regime. Significance threshold:
$p < 0.00033$ (Bonferroni-corrected for 30 factor pairs).}
\label{tab:main}
\small
\begin{tabular}{llccc}
\toprule
Regime & Direction & $p$-value & Lag & Significant \\
\midrule
Normal & HML $\to$ SMB & $\mathbf{8.75 \times 10^{-9}}$ & 1 & \textbf{Yes} \\
Normal & SMB $\to$ HML & $0.864$ & 1 & No \\
\midrule
Elevated & HML $\to$ SMB & $0.004$ & 1 & Suggestive$^*$ \\
Elevated & SMB $\to$ HML & $0.897$ & 1 & No \\
\midrule
Crisis & HML $\to$ SMB & $0.695$ & 1 & No \\
Crisis & SMB $\to$ HML & $0.294$ & 1 & No \\
\bottomrule
\multicolumn{5}{l}{\footnotesize $^*$Below 0.01 but above Bonferroni threshold $\alpha/30 = 0.00033$.}
\end{tabular}
\end{table}

\textbf{Key finding under confirmatory threshold:} only Normal-regime
HML$\to$SMB survives Bonferroni correction.
\begin{itemize}
    \item \textbf{Normal (confirmatory):} Strong lag-1 effect
    ($p = 8.75 \times 10^{-9}$), robust to HAC correction.
    \item \textbf{Elevated (suggestive):} HML$\to$SMB is directional at
    conventional levels ($p = 0.004$) but does not pass the paper's
    multiple-testing threshold.
    \item \textbf{Crisis (null):} No detectable lag-1 Granger effect in either
    direction.
\end{itemize}

\textbf{Robustness to heteroskedasticity.}
Financial returns exhibit volatility clustering, which may inflate standard Granger
$F$-test significance. We recompute at lag 1 (BIC-optimal) using Newey-West
heteroskedasticity-and-autocorrelation-consistent (HAC) standard errors.
Normal-regime HML$\to$SMB survives HAC correction
($p = 9.45 \times 10^{-8}$ HAC vs.\ $p = 8.75 \times 10^{-9}$ standard).
Elevated HML$\to$SMB weakens under HAC ($p = 0.034$), and Crisis remains
non-significant ($p = 0.711$ HAC vs.\ $p = 0.695$ standard).

\textbf{Incremental $R^2$.} Table~\ref{tab:r2} quantifies the economic magnitude. Adding
HML lags (1--15) to the SMB autoregression increases $R^2$ by 1.04\% in Normal,
1.32\% in Elevated, and 0.56\% in Crisis. Normal and Elevated pass conventional
significance levels ($p = 5.98 \times 10^{-5}$ and $p = 0.002$), while Crisis
does not ($p = 0.988$). In the reverse direction, SMB lags add 1.35\% in Crisis
with $p = 0.562$, so we do not claim directional confirmation there.

\begin{table}[t]
\centering
\caption{Incremental $R^2$ from adding HML lags (1--15) to SMB autoregression by regime.
$\Delta R^2$ = $R^2_{\text{unrestricted}} - R^2_{\text{restricted}}$.}
\label{tab:r2}
\small
\begin{tabular}{lcccc}
\toprule
Regime & $n$ & $R^2_{\text{AR}}$ & $\Delta R^2$ & $p$-value \\
\midrule
Normal & 4,361 & 0.61\% & 1.04\% & $5.98 \times 10^{-5}$ \\
Elevated & 2,690 & 0.59\% & 1.32\% & $0.002$ \\
Crisis & 981 & 1.01\% & 0.56\% & $0.988$ \\
\bottomrule
\end{tabular}
\end{table}

\begin{figure}[t]
\centering
\includegraphics[width=\columnwidth]{figures/granger_heatmap.pdf}
\caption{All-pairs Granger causality by regime ($-\log_{10}(p)$, lag=5).
Bordered cells with asterisks survive Bonferroni correction.
The strongest regime-differential pair is MKT$\to$SMB; HML$\to$SMB appears as
a secondary differential pattern.}
\label{fig:heatmap}
\end{figure}

\textbf{All-pairs perspective.} Figure~\ref{fig:heatmap} shows Granger causality
$p$-values for all directed factor pairs across three regimes. Several pairs
show regime-dependent variation; MKT$\to$SMB has the largest differential.
HML$\to$SMB ranks fourth in this differential metric, with strongest evidence in
Normal ($p = 1.0 \times 10^{-6}$ at lag 5), suggestive Elevated evidence
($p = 0.013$), and no Crisis evidence ($p = 0.880$). We therefore treat
HML$\to$SMB as one regime-sensitive relation among several, not the dominant pair.

\textbf{Lag specification and persistent structure.} We tested lags 1--15 for
HML$\to$SMB across regimes (Figure~\ref{fig:lag_sensitivity}). In Normal, the
effect is strongest at lag 1 and remains significant through lag 15
($p = 5.98 \times 10^{-5}$ at lag 15). Elevated shows weak-to-moderate
evidence across lags (e.g., lag 15 $p = 0.002$), while Crisis is consistently
null (lag 1 $p = 0.695$, lag 15 $p = 0.988$). We present Elevated/Crisis lag
patterns as exploratory because they do not meet the confirmatory threshold.

\subsection{Hypothesized Economic Mechanism}
\label{sec:interpretation}

We offer the following interpretation as a \emph{hypothesis}, not a verified mechanism:

\textbf{Normal-regime lag-1 link (confirmatory):} During calmer periods,
cross-factor inventory management and short-horizon rebalancing may transmit
value shocks into size returns with a one-day delay.

\textbf{Stress-channel hypothesis (exploratory):} Portfolio-overlap evidence in
FF25 remains consistent with a deleveraging channel in subsets of stress data
($\rho_s = 0.35$, permutation $p = 0.046$), but primary Crisis-regime
factor-level Granger tests are not significant in the leakage-safe fit.

\textbf{Alternative explanations we cannot rule out:}
\begin{itemize}
    \item A common driver (e.g., funding liquidity) affects both factors with different lags
    \item Information about distress propagates through markets, reaching Value-sensitive
    investors before Size-sensitive investors
    \item Statistical artifact of regime-conditional estimation
\end{itemize}

Greenwood and Thesmar~\cite{greenwood2011stock} document that stocks with
concentrated institutional ownership exhibit correlated fragility during
liquidation events---consistent with our hypothesized deleveraging channel.
We provide exploratory portfolio-level evidence for this overlap in
Section~\ref{sec:overlap_evidence}.

\subsection{Portfolio Overlap Evidence}
\label{sec:overlap_evidence}

The deleveraging hypothesis predicts that the HML$\to$SMB Granger link should
be strongest for portfolios that load on \emph{both} factors simultaneously.
The Small-HighBM portfolio sits in both the SMB long leg (small stocks) and
the HML long leg (high book-to-market), making it the primary candidate
transmission channel.

We test this prediction using the Fama-French 25 Size$\times$BM portfolios
(daily, 1990--2024). For each of the 25 portfolios, we run a Granger test
of HML$\to$Portfolio at lag 1 (BIC-selected in Table~\ref{tab:main}) within
the Crisis regime.

\textbf{Spatial gradient test.} We define an ``overlap score'' for each
portfolio as $(4 - \text{size quintile}) + \text{BM quintile}$, ranging from
0 (Big/LowBM, no overlap) to 8 (Small/HighBM, maximal overlap). The
confirmatory test is a Spearman rank correlation between overlap score and
$-\log_{10}(p)$ of the Granger test across 25 portfolios, with significance
assessed by spatial permutation (10{,}000 shuffles of grid assignments;
time series are never touched).

\textbf{Result:} $\rho_s = 0.35$, permutation $p = 0.046$
(Figure~\ref{fig:ff25}). Granger significance concentrates in small-cap
portfolios: all five Small quintile portfolios are nominally significant
($p < 0.01$), with three surviving Bonferroni correction ($\alpha/25 = 0.002$):
S/L ($p = 2.5 \times 10^{-6}$, $\Delta R^2 = 1.47\%$),
S/3 ($p = 0.0014$), and S/H ($p \approx 0$, $\Delta R^2 = 1.79\%$).
One Big-quintile portfolio is nominally significant ($p < 0.01$), but none
survive Bonferroni. This spatial pattern is inconsistent
with a uniform factor-level effect and supports the overlap channel.

\textbf{Overlap fraction.} At lag 15 (matching Table~\ref{tab:r2}),
$\Delta R^2(\text{HML}\to\text{S/H}) / \Delta R^2(\text{HML}\to\text{SMB})
= 0.392$: the Small-HighBM portfolio alone accounts for 39\% of the
HML$\to$SMB incremental $R^2$.

\begin{figure}[t]
\centering
\includegraphics[width=\columnwidth]{figures/ff25_overlap_granger_heatmap.pdf}
\caption{FF 25 Portfolio Overlap Analysis. Panel A: $-\log_{10}(p)$ for
HML$\to$Portfolio Granger tests (Crisis, lag=1). Significance concentrates in
Small quintile; $^{**}$Bonferroni-significant. Panel B: Spatial gradient---overlap
score vs.\ Granger significance ($\rho_s = 0.35$, permutation $p = 0.046$).}
\label{fig:ff25}
\end{figure}

\subsection{Early Warning Performance}
\label{sec:early_warning}

The Student-$t$ HMM detects crisis regimes before market peaks (Table~\ref{tab:warning}).
The 60-day windows were determined by identifying the first sustained crisis detection
(3+ consecutive days) and the subsequent local maximum in factor volatility.

\begin{table}[t]
\centering
\caption{Early Warning Lead Time. First Detection = first day of 3+ consecutive
Crisis regime assignments.}
\label{tab:warning}
\scriptsize
\begin{tabular}{@{}lccc@{}}
\toprule
Event & First Det. & Vol. Peak & Lead \\
\midrule
Lehman 2008 & May 9, 2008 & Sep 15, 2008 & 129 days \\
EU Crisis 2011 & None & Aug 8, 2011 & None \\
COVID 2020 & None & Mar 23, 2020 & None \\
\bottomrule
\end{tabular}
\end{table}

\section{Conclusion}

The recomputed, leakage-safe primary fit supports a narrower conclusion than
earlier drafts. HML$\to$SMB is strongly significant in the Normal regime at lag 1
($p = 8.75 \times 10^{-9}$; HAC $p = 9.45 \times 10^{-8}$). Elevated evidence is
suggestive but non-confirmatory under the paper's multiple-testing threshold,
and Crisis evidence is null at lag 1 ($p = 0.695$). The frozen out-of-sample
aggregate test is positive at conventional levels ($p = 0.041$), but event-level
results remain mixed.

\textbf{Does establish:}
\begin{itemize}
    \item Confirmatory Normal-regime HML$\to$SMB predictive precedence
    \item Exploratory FF25 overlap evidence ($\rho_s = 0.35$, permutation $p = 0.046$)
    \item No evidence of tradable alpha (Sharpe $= -0.07$)
    \item Practical utility of a hybrid detector for VaR calibration
    (5.60\% violations, CC $p = 0.336$)
\end{itemize}

\textbf{Does not establish:}
\begin{itemize}
    \item Uniform significance across every regime
    \item Uniform event-level support across all stress episodes
    \item Structural causality (common drivers could explain the pattern)
    \item Trading profitability
\end{itemize}

\textbf{Practical value:} A hybrid VaR model combining HMM regime detection
with a realized-volatility fallback achieves the closest coverage to the 5\%
target (5.60\%, $p_{\text{CC}} = 0.336$), detecting 47.8\% of COVID-2020
where the HMM alone detected 0\%. Algorithm~\ref{alg:monitor} formalizes
the monitoring protocol with all parameters calibrated on training data.

For practitioners, the key insight is that factor lead-lag relationships are not
merely stronger or weaker across regimes. In the recomputed primary fit,
the clearest actionable signal is Normal-regime lag-1 HML$\to$SMB. Crisis-regime
monitoring is better treated as stress-sensitivity analysis than as a
confirmatory directional signal. Real-time implementation still benefits from
the volatility override: the frozen HMM alone classified 0\% of COVID-2020 as
Crisis, while the hybrid detector reached 47.8\% detection with one
pre-specified threshold parameter.

\bibliography{references}
\end{document}
